\documentclass[conference]{IEEEtran}
\IEEEoverridecommandlockouts
\usepackage{cite}
\usepackage{amsmath,amssymb,amsfonts}
\usepackage{algorithm}
\usepackage{algorithmic}
\usepackage{graphicx}
\usepackage{textcomp}
\usepackage{xcolor}
\usepackage{booktabs}
\usepackage{microtype}
\usepackage[hidelinks]{hyperref}

\def\BibTeX{{\rm B\kern-.05em{\sc i\kern-.025em b}\kern-.08em
    T\kern-.1667em\lower.7ex\hbox{E}\kern-.125emX}}
\begin{document}

\title{A Bounded-Latency Virtual Mentor for Introductory Programming Feedback}

\author{
    \IEEEauthorblockN{Alimzhan Koshikov}
    \IEEEauthorblockA{
        \textit{Master Degree Student, School of Software Engineering}\\
        Astana IT University (AITU)\\
        Astana, Kazakhstan\\
        225148@astanait.edu.kz\\
        \href{https://orcid.org/0009-0003-3234-6487}{0009-0003-3234-6487}
    }
    \and
    \IEEEauthorblockN{Tutkyshbayeva Shyryn}
    \IEEEauthorblockA{
        \textit{PhD, School of Software Engineering}\\
        Astana IT University (AITU)\\
        Astana, Kazakhstan\\
        sh.tutkyshbayeva@astanait.edu.kz\\
        \href{https://orcid.org/0000-0002-1419-2571}{0000-0002-1419-2571}
    }
}

\maketitle

\begin{abstract}
Introductory programming remains challenging for many CS1 students because they must learn syntax, problem decomposition, and basic program reasoning at the same time. In large classes, timely individualized feedback is difficult to provide consistently, while existing automated feedback is often delayed, generic, or limited to correctness checks. This paper presents a bounded-latency virtual mentor for introductory programming that delivers real-time, pedagogically constrained feedback during active problem solving.

The proposed system combines continuous code analysis, a lightweight student state, and a policy layer that selects an appropriate intervention type, such as a guiding question, a conceptual hint, a code highlight, or no feedback. Unlike unconstrained conversational assistants, the system is designed to regulate feedback timing and granularity in order to support student independence rather than replace problem solving.
\end{abstract}

\begin{IEEEkeywords}
Virtual mentor, CS1, introductory programming, programming education, real-time feedback, student modeling, pedagogical constraints
\end{IEEEkeywords}

\section{Introduction}

Introductory programming (CS1) remains difficult for first-year students because they must simultaneously learn syntax, develop semantic understanding, and build problem-solving habits. In large classes, beginners often repeat the same compiler and logic errors, while instructors and teaching assistants cannot always provide immediate and individualized support at scale \cite{DesirableCharacteristicsforAITeachingAssistantsinProgrammingEducation,FeasibilityStudyofAugmentingTeachingAssistantswithAI}.

Automated assessment improves scalability through compiler messages, unit tests, and static checks, but such feedback is often binary or generic and does not explicitly address misconceptions or the level of support a student actually needs. Prior work in computing education shows that feedback quality depends not only on correctness, but also on timing, pedagogical framing, and whether the intervention preserves student autonomy rather than simply producing a working answer \cite{SystematicLiteratureReviewonLLM,DesirableCharacteristicsforAITeachingAssistantsinProgrammingEducation}.

Recent AI-based teaching assistants have shown that near real-time support is feasible in large-scale educational settings. For example, the CS50 AI assistant reported high adoption and low operating cost at MOOC scale \cite{ImprovingAIinCS50}. At the same time, prior studies also highlight important risks: students may overuse assistants for immediate help, explanations may be inaccurate, and higher perceived usefulness does not necessarily translate into better learning outcomes \cite{FeasibilityStudyofAugmentingTeachingAssistantswithAI,PatternsofLLMPoweredProgrammingAssistant,SystematicLiteratureReviewonLLM}. These findings suggest that educational assistants should not only generate responses, but should also regulate \emph{when} help is given, \emph{what kind} of help is given, and \emph{how much} information is disclosed \cite{CodeAid,IRIS,DesirableCharacteristicsforAITeachingAssistantsinProgrammingEducation}.

This paper presents a virtual mentor for CS1 that is designed as a bounded real-time decision pipeline rather than as a general conversational assistant. The system combines continuous code analysis, a lightweight student state, and a pedagogical policy layer that selects an intervention type under latency constraints. If a language model is used, it is treated only as a realization component for phrasing feedback, not as the core decision-making mechanism.

The main contribution of this work is not another generic AI coding assistant, but a prototype architecture for \emph{policy-guided, real-time mentoring} during active programming. In contrast to systems that respond only after submission or produce unconstrained natural-language help, the proposed mentor explicitly separates analyzer signals, student state tracking, action selection, and feedback realization. This separation improves interpretability, supports bounded response time, and makes the system suitable for future comparison between rule-based and learned policy variants.

The main contributions of this work are as follows:
\begin{itemize}
\item a modular virtual mentor architecture that separates real-time code analysis, lightweight student-state tracking, pedagogical action selection, and feedback realization;
\item a policy-guided feedback mechanism that regulates timing and granularity of help while avoiding direct solution disclosure;
\item a pilot prototype evaluation based on event-level interaction logs, demonstrating bounded-latency operation and preliminary action-level effectiveness trends.
\end{itemize}

\section{Related Work}

Research on AI support for programming education commonly falls into three broad directions: intelligent tutoring systems, automated feedback generation, and LLM-based assistants. The proposed virtual mentor draws from all three, but combines them with a different emphasis on bounded-latency interaction and policy-level control.

\subsection*{A. Intelligent Tutoring Systems for Programming}

Intelligent tutoring systems (ITS) personalize support by modeling learner state and adapting feedback or task sequencing. Recent systems increasingly aim to augment instructors rather than fully automate instruction. ProgTutor, for example, combines adaptive sequencing with automated evaluation while keeping instructors in the loop \cite{ProgTutor}. A common limitation, however, is that many ITS intervene at discrete checkpoints or after submission, which reduces their usefulness during the moment-to-moment process of code construction.

The proposed virtual mentor targets this in-process setting by monitoring code continuously and selecting interventions during active problem solving rather than only after a completed attempt.

\subsection*{B. Automated Feedback Generation Systems}

Automated feedback is widely used in CS1 because it scales well and is easy to integrate into programming courses. Typical implementations rely on unit tests, static analysis, or compiler messages. These methods are useful for fast validation, but they often remain shallow: they say whether something is wrong without adapting the explanation to the learner's misunderstanding. Kumar et al. note that many automated feedback systems emphasize syntax correction or output correctness rather than concept-level guidance, while reproducibility and personalization remain open challenges \cite{NavigatingAutomatedFeedback}.

This limitation motivates the combination used in our work: code analysis alone is not sufficient, so it is coupled with lightweight student modeling and a pedagogical decision layer.

\subsection*{C. LLM-Based Programming Assistants}

Compared with rule-based systems, LLM-based assistants are more flexible in generating explanations and hints in natural language \cite{SystematicLiteratureReviewonLLM}. Systems such as CodeAid and Iris show that LLM-supported assistance can be valuable when pedagogical constraints are built into the interaction design \cite{CodeAid,IRIS}. At the same time, classroom evidence highlights important risks. Students may seek immediate answers instead of engaging with the underlying concept, and the assistant may still produce explanations that are inaccurate, overly strong, or poorly aligned with course objectives \cite{PatternsofLLMPoweredProgrammingAssistant,FeasibilityStudyofAugmentingTeachingAssistantswithAI}. At larger scale, CS50 also emphasizes the need for continuous human feedback and systematic evaluation to keep AI support pedagogically aligned \cite{ImprovingAIinCS50}.

\subsection*{D. Research Gap and Positioning}

Existing work does not consistently combine three requirements in a single CS1 support system: (i) intervention during active code writing rather than only after submission, (ii) an explicit lightweight student state that captures repeated difficulties, and (iii) a pedagogical policy that regulates feedback timing, type, and disclosure level. Intelligent tutoring systems provide personalization, but they often rely on heavier domain modeling or operate at coarser interaction granularity. LLM-based assistants offer fluent explanations, but many of them function as open-ended help tools without explicit student-state tracking or controlled intervention logic \cite{ProgTutor,NavigatingAutomatedFeedback,CodeAid,IRIS}.

The virtual mentor proposed in this work is positioned between these two lines of research. It does not attempt to replace instructors, nor does it rely on unconstrained response generation. Instead, it treats mentoring as a bounded-latency decision process: analyzer signals are combined with a lightweight student state, a policy selects an action, and the selected action is then realized as feedback. This design shifts the technical focus from generic answer generation to controlled pedagogical intervention.

Table~\ref{tab:feedback_comparison} summarizes representative systems and highlights the positioning of the proposed virtual mentor relative to prior approaches.

\begin{table*}[t]
\centering
\caption{Positioning of the Proposed Virtual Mentor Against Representative Programming Feedback Systems}
\label{tab:feedback_comparison}
\renewcommand{\arraystretch}{1.2}
\begin{tabular}{p{2.5cm} p{3.5cm} p{4.8cm} p{3.9cm}}
\hline
\textbf{System} &
\textbf{Core Technology} &
\textbf{Main Feedback Mode} &
\textbf{Reported Focus / Metric} \\
\hline

AIF~3.0 \cite{AdaptiveImmediateFeedbackforBlockBasedProgramming} &
Hybrid data-driven + expert-based &
Progress indicators and subgoal-oriented feedback &
Subgoal detection accuracy: 88\% \\

CodeAid \cite{CodeAid} &
GPT-3.5-Turbo &
Pseudocode-oriented explanations and guided support &
Technical correctness: 79\% \\

Verifix \cite{NavigatingAutomatedFeedback} &
Program analysis &
Verified repair suggestions &
Feedback quality score: 3.8/5.0 \\

DeepFix \cite{NavigatingAutomatedFeedback} &
Neural error-correction model &
Syntax error correction &
Full error correction rate: 27\% \\

Proposed Virtual Mentor &
Real-time analyzer + student state + policy layer + optional LLM realization &
Pedagogically constrained in-process feedback during code writing &
Pilot prototype results reported in Sec.~\ref{sec:evaluation} \\

\hline
\end{tabular}
\end{table*}

\section{Virtual Mentor System Architecture}

The proposed virtual mentor is a bounded real-time decision pipeline that maps a stream of student code states to pedagogically constrained feedback actions. The architecture separates code analysis, student-state tracking, policy selection, and feedback realization in order to improve modularity, interpretability, and latency control.

Unlike standalone LLM assistants that act as external help tools, the mentor is embedded into the programming workflow and observes code changes incrementally. If natural-language generation is used, it operates only at the realization stage. The core decision logic remains outside the language model and is driven by analyzer signals and student-state features.

\subsection*{A. Data Flow}

As shown in Fig.~\ref{fig:architecture}, code edits are processed by an analyzer that performs syntactic checks, lightweight structural analysis, and detection of common novice error patterns. The analyzer output is passed to a profile manager that maintains a compact student state, including repeated error categories and recent feedback history. A policy layer then selects the appropriate feedback action, and the selected action is realized as a concrete intervention, such as a guiding question, a conceptual hint, or a code-region highlight.

This separation is important for two reasons. First, it allows the system to remain responsive under real-time constraints by keeping the expensive components outside the critical decision path. Second, it makes the intervention logic explicit, which is useful both for pedagogical control and for future comparison between rule-based and learned policy variants.

\subsection*{B. Formal System Model}

Let $c_t$ denote the student code state after an edit at time $t$. The analyzer
\[
A: c_t \rightarrow e_t
\]
maps $c_t$ to an error descriptor $e_t \in \mathcal{E}$, where $\mathcal{E}$ is a finite set of novice-relevant error categories.

The student state is represented as a lightweight vector
\[
s_t = (n_{e_1}, n_{e_2}, \dots, n_{e_k}, h_t),
\]
where $n_{e_i}$ is the cumulative count of error type $e_i$ observed up to time $t$, and $h_t$ denotes compact history features such as the most recent intervention type.

Given $(e_t, s_t)$, the policy selects an action
\[
\pi: (e_t, s_t) \rightarrow a_t,
\]
where $a_t \in \mathcal{A}$ and
\[
\mathcal{A}=\{\textsc{GQ}, \textsc{CH}, \textsc{HL}, \textsc{NF}\}
\]
corresponds to \emph{Guiding Question}, \emph{Conceptual Hint}, \emph{Code Highlight}, and \emph{No Feedback}.

In the current prototype, $\pi$ is implemented as a rule-based policy over analyzer outputs and student-state features. The same interface also supports a learned policy extension, where action selection can be formulated as supervised multi-class classification:
\[
\hat{a} = \arg\max_{a \in \mathcal{A}} P(a \mid x),
\]
with feature vector $x$ constructed from the current error category, recurrence counts, and recent interaction context.

\subsection*{C. Core Components}

\textbf{Analyzer:} Processes incoming code states and extracts pedagogically relevant signals, including compile status, novice error categories, and suspicious code regions.

\textbf{Profile Manager:} Maintains a lightweight and interpretable student state with low computational overhead. The design favors responsiveness and transparency over heavy knowledge-tracing machinery.

\textbf{Policy Layer:} Chooses whether and how the system should intervene. This component is distinct from code analysis: it does not detect errors itself, but decides how analyzer signals should be translated into pedagogical action.

\textbf{Feedback Realization:} Produces the final response corresponding to the chosen action. This may be implemented through templates or a language model, but in both cases the output is constrained by the selected action type and the no-solution-disclosure rule.

\section{Real-Time Feedback Mechanism}

The feedback mechanism is designed for in-process support during active problem solving rather than post-submission correction. Prior work on LLM-based programming assistants shows that immediate help is useful, but overly detailed or poorly timed assistance may encourage trial-and-error behavior and reduce conceptual engagement \cite{PatternsofLLMPoweredProgrammingAssistant}. For this reason, the proposed system does not treat every detected issue as a trigger for the same kind of response.

\subsection*{A. Feedback Generation}

\begin{figure}[t]
\centering
\includegraphics[width=\linewidth]{arch-SIST.png}
\caption{High-level architecture of the virtual mentor system.}
\label{fig:architecture}
\end{figure}

When a learner modifies code, the analyzer first produces a compact description of the current state, including error-related signals and possible suspicious regions. These signals are recorded in the student state, allowing the system to distinguish between new difficulties and repeated ones. The policy layer then selects one of four intervention types: a guiding question, a conceptual hint, a code highlight, or no feedback.

This design reflects a pedagogical assumption that not all support should be equally explicit. In some cases, highlighting a suspicious region is sufficient. In others, repeated difficulties may justify a conceptual hint. A guiding question is used when the system aims to prompt reasoning without disclosing the next implementation step directly. Complete solutions are intentionally avoided.

\subsection*{B. Timing Constraints}

To preserve interactivity, the total response time is constrained by
\[
T_{total} = T_{analysis} + T_{policy} + T_{feedback},
\]
where $T_{analysis}$ denotes analyzer processing time, $T_{policy}$ denotes action-selection time, and $T_{feedback}$ denotes feedback realization time.

The prototype is designed so that the policy stage remains lightweight and does not dominate the latency budget. In the pilot interaction logs, the end-to-end latency remained within interactive bounds, with mean $T_{total}=\textbf{741}$\,ms and P95 $\textbf{809}$\,ms. The mean component times were $\textbf{618}$\,ms for analysis, $\textbf{26}$\,ms for policy selection, and $\textbf{12}$\,ms for feedback realization, with the remaining delay attributable to platform overhead.

\section{Experimental Evaluation}\label{sec:evaluation}

\subsection*{A. Evaluation Artifact}

The current evaluation is based on a pilot event-level interaction log produced by the prototype rather than on a completed classroom-scale comparative study. The artifact contains \textbf{60} logged interaction events collected during representative CS1 problem-solving sessions. Each event records analyzer output, selected feedback action, latency measurements, and whether the issue was resolved after the intervention.

This pilot evaluation is intended to assess prototype feasibility, responsiveness, and preliminary action-level behavior. It does not yet establish longitudinal learning gains, statistical classroom effects, or superiority over external benchmark systems.

Each logged interaction is represented as
\[
x_i = (e_i, c_i, t^{analysis}_i, t^{policy}_i, t^{feedback}_i, a_i, s_i),
\]
where $e_i$ denotes the current error category, $c_i$ the recurrence count, $t^{analysis}_i$, $t^{policy}_i$, and $t^{feedback}_i$ the measured latency components, $a_i$ the selected feedback action, and $s_i$ the event outcome indicator.

\subsection*{B. Metrics}

Two groups of metrics were considered. First, responsiveness was measured through end-to-end latency and its component breakdown:
\[
\bar{T} = \frac{1}{N} \sum_{i=1}^{N} T_{total}^{(i)}.
\]

Second, we used event-level effectiveness proxies derived from the interaction logs, including event success rate and time-to-fix. In this context, an event is counted as successful if the issue associated with the intervention was resolved after the selected feedback action. These metrics provide an operational view of the prototype, but they should not be interpreted as full learning-outcome measures.

\subsection*{C. Results}

Table~\ref{tab:results} summarizes the pilot results. The main purpose of these results is to show that the proposed architecture can operate within interactive latency constraints while producing different intervention types in a functioning mentoring loop.

\begin{table}[t]
\centering
\caption{Pilot Results from Event-Level Interaction Logs}
\label{tab:results}
\renewcommand{\arraystretch}{1.15}
\begin{tabular}{lc}
\hline
\textbf{Metric} & \textbf{Value} \\
\hline
Events (N) & 60 \\
Mean latency $\bar{T}$ (ms) & 741 \\
P95 latency (ms) & 809 \\
Mean analysis/policy/feedback (ms) & 618 / 26 / 12 \\
Platform overhead (ms) & 85 \\
Event success rate (\%) & 66.7 \\
Median time-to-fix (ms) & 2450 \\
P95 time-to-fix (ms) & 3595 \\
\hline
\end{tabular}
\end{table}

\begin{table}[t]
\centering
\caption{Event Success Rate by Feedback Action}
\label{tab:action_success}
\renewcommand{\arraystretch}{1.15}
\begin{tabular}{lcc}
\hline
\textbf{Action} & \textbf{N} & \textbf{Success rate} \\
\hline
CODE\_HIGHLIGHT & 23 & 1.00 \\
CONCEPTUAL\_HINT & 15 & 1.00 \\
GUIDING\_QUESTION & 19 & 0.25 \\
NO\_FEEDBACK & 3 & 0.00 \\
\hline
\end{tabular}
\end{table}

\section{Discussion}

\subsection*{A. Interpretation of Results}

The pilot results suggest that the proposed mentoring loop can operate under interactive time constraints. With a mean end-to-end latency of \textbf{741}\,ms and P95 of \textbf{809}\,ms, the prototype remained within a response range that is compatible with in-process feedback during code writing. This is important because the main technical goal of the system is not only to generate help, but to do so without disrupting the programming flow.

The action-level breakdown also provides an initial indication that different intervention types may have different short-term effects. In the current trace, code highlighting and conceptual hints were associated with higher immediate resolution rates than guiding questions. This does not yet justify strong causal claims, but it does support the idea that feedback type matters and that action selection should be treated as an explicit design problem rather than as a by-product of response generation.

More broadly, the pilot supports the feasibility of the architectural separation used in this work. The analyzer, student state, policy layer, and feedback realization operated as distinct components inside one real-time loop, which is the main systems contribution of the prototype.

\subsection*{B. Limitations}

This study has several limitations. First, the present evaluation is based on a pilot interaction log with a limited number of events and should therefore be interpreted as prototype-level evidence rather than classroom-scale validation. Second, the current artifact does not provide inferential comparison against external benchmark systems or against a completed control-group study. Third, the reported effectiveness measures are short-term event-level proxies and do not establish learning gains, retention, or transfer. Finally, the student state remains intentionally lightweight; richer learner modeling may improve adaptivity, but would also increase implementation complexity and data requirements.

\subsection*{C. Practical Implications}

Despite these limitations, the results indicate that real-time personalized support can be introduced as a middle layer between static automated checks and human instruction. In this role, the mentor does not replace instructors or teaching assistants. Instead, it helps filter routine support situations, provides bounded assistance during problem solving, and preserves instructional intent through explicit control of intervention type.

\subsection*{D. Ethical Considerations}

The system is designed to support student autonomy by avoiding direct solution disclosure and prioritizing formative guidance. Interaction logs are stored for system analysis and should be anonymized in real deployments. Since AI-assisted help may encourage over-reliance if used without safeguards, the policy layer is intended to regulate feedback timing and explicitness rather than maximize immediate answer delivery.

\section{Future Work}

Several directions remain for future work. First, the prototype should be evaluated on a larger and more diverse set of CS1 tasks, including comparative baselines such as non-personalized automated feedback and fixed rule-only intervention policies. Second, the current rule-based action selector can be extended with a learned policy model trained on labeled interaction cases while preserving the same feature interface. Third, broader classroom studies are needed to test whether short-term event-level effectiveness translates into measurable learning outcomes, retention, and transfer. Finally, the student-state representation can be extended beyond simple recurrence counts to include temporal learning dynamics and teacher-in-the-loop adjustment of policy behavior.

\section{Conclusion}

This paper presented a bounded-latency virtual mentor for introductory programming that combines real-time code analysis, lightweight student-state tracking, pedagogical action selection, and constrained feedback realization.

The pilot prototype evaluation showed that the system can operate within interactive latency limits, with mean end-to-end latency of \textbf{741}\,ms and P95 of \textbf{809}\,ms. Across \textbf{60} event-level interactions, the prototype achieved an event success rate of \textbf{66.7\%} and a median time-to-fix of \textbf{2.45}\,s. These results should be interpreted as feasibility evidence rather than as final learning-outcome validation, but they support the practical viability of policy-guided, real-time mentoring for CS1.

The main contribution of the work lies in the architectural separation of analyzer signals, student-state tracking, policy-based intervention, and feedback realization under bounded latency. This provides a concrete foundation for future comparative and classroom-scale studies.

\section*{Acknowledgment}

This work was carried out as part of the master’s degree program at Astana IT University. The authors thank Astana IT University for providing the academic environment and computing infrastructure that supported the development of the prototype.

\nocite{Francisco2022}
\nocite{Chrysafiadi2022}
\nocite{pankiewicz2023largelanguagemodelsgpt}
\nocite{AutomatedAssessmentToolsForGrading2022}
\nocite{AdaptiveImmediateFeedbackforBlockBasedProgramming}
\nocite{FeedbackSystemSupportingStudents2022}
\nocite{EvaluationofLLMtoolsforfeedbackgeneration}
\nocite{AutomatingHumanTutorStyleProgrammingFeedback}
\nocite{AITutoringinSDEducation}
\nocite{AutomatedGradingandFeedbackToolsforProgrammingEducationSystematicReview}
\bibliographystyle{IEEEtran}
\bibliography{references}

\end{document}
